\documentclass[12pt, a4paper]{article}
\title{Sprawozdanie z projektu, Metody modelowania matematycznego}
\author{Joanna Jabczyńska, s198075}
\date{May, 2025}

\usepackage{graphicx}
\usepackage{tcolorbox}
\usepackage[polish]{babel}
\usepackage[utf8]{inputenc}
\usepackage[T1]{fontenc}
\usepackage{lmodern}
% --- \usepackage{fancyhdr}
% --- \pagestyle{fancy}

\begin{document}
\maketitle
% --- \thispagestyle{fancy}
\section* {Polecenie - zadanie 4.}
Dany jest układ mechaniczny przedstawiony na poniższym rysunku:

\includegraphics[width=\textwidth]{image}
Należy wyprowadzić model układu oraz zaimplementować go w symulacji. Symulator
powinien umożliwiać pobudzenie układu przynajmniej trzema rodzajami synagłów
wejściowych (prostokąny o skończonym czasie trwania, trójkątny, harmoniczny). Symulator
powinien umożliwiać zmianę wszystkich parametrów układu oraz sygnałów wejściowych.
Należy użyć metody Rungego-Kutty 4-go rzędu oraz metody Eulera oraz na wspólnym
wykresie pokazać wyniki symulacji (prędkości i położenia wału J2) z obu tych metod.

\section*{Założenia projektu}
\subsection{Potrzebne obliczenia}

\[J_e = J_1{\cdot}{(\frac{n_1}{n_2})}^2+J_2\]
\[k_e = k{\cdot}{(\frac{n_1}{n_2})}^2\]


\[J_e{\cdot}\ddot{\phi_2}+b{\cdot}\dot{\phi_2}+k_e{\cdot}{\phi_2}=T{\cdot}\frac{n_2}{n_1}\]

\[x_1 = \phi_2\]
\[\dot{x_1} = x_2\]
\[\dot{x_2} = \ddot{\phi_2}\]

\[J_e{\cdot}\dot{x_2}+b{\cdot}x_2+k_ex_1=T{\cdot}\frac{n_2}{n_1}\]
\[\dot{x_2} = \frac{T{\cdot}\frac{n_2}{n_1}-b{\cdot}x_2-k_e{\cdot}x_1}{J_e}\]

\subsection{Potrzebne elementy kodu}
	\begin{itemize}
		\item Funkcje implemetujące metody numeryczne: Eulera i Rundego-Kutty 4-go rzędu
		\item Wizualizacja wyników
		\item Interakcja użytkownika z parametrami obiektu i sygnału wejściowego
		\item Pobieranie danych wyjściowych przez użytkownika w pliku .csv
	\end{itemize}
\section*{Realizacja techniczna}
	\begin{enumerate}
		\item Wykorzystanie języka Python ze względu na jego prostotę
		\item Skorzystanie z bibliotek:
			\begin{itemize}
				\item NumPy - operacje i stałe matematyczne oraz wartości funkcji sinus
				\item Streamlit - wizualizacja danych wyjściowych oraz stworzenie interfejsu graficznego użytkownika
				\item Pandas - przechowywanie i organizacja danych obliczeniowych
				\item Zipfile - pakowanie plików .csv w plik .zip, aby umożliwić pobranie przez użytkownika
			\end{itemize}
	\end{enumerate}
\section*{Wnioski projektu}
	\begin{itemize}
		\item Obie metody dają bliskozbieżne wyniki dla różniczkowania $\dot{\phi_2}$ (prędkość wału), jednak potrafi wskazywać rozbieżność o stałą dla położenia wału.
		\item Metoda Eulera wymaga mniejszej liczby operacji, jest zatem bardziej oszczędna
	\end{itemize}
\subsection{Wpływ parametrów układu}
	\begin{itemize}
		\item \emph{b} (tłumienie): Wygaszenie tłumienia (wartość parametru równa 0) 
		\item \emph{b} Wziększenie tłumienia powoduje stabilizowanie się sygnałów wyjściowych wokół określonej wartości. 
		\item \emph{k} (sprężystość): Wziększenie współczynnika sprężystości powoduje zwiększenie się amplitudy elementów harmonicznych 
		\item \emph{J} (bezwładność): Zwiększenie bezwładności powoduje spowolnienie układu (zwiększenie się stałej czasowej)
		\item \emph{n} Zwiększenie stosunku $\frac{n_1}{n_2}$ powoduje zanikanie charakterystyki harmonicznej, zaś zwiększenie wzmaga oscylacje i spowalnia ich zanikanie
	\end{itemize}
\subsection{Wpływ parametrów sygnału}
	\begin{itemize}
		\item Zwiększenie częstotliwości sygnału powoduje wypłaszczenie się krzywej położenia wału
		\item Zwiększenie amplitudy powoduje wzrost wartości maksymalnej
	\end{itemize}
\end{document}